\section{Cayley 圆盘化与 Jensen 缺陷：无参紧化的 RH 等价式}\label{app:cayley-disk-jensen}
本附录补入一条与正文 \S\ref{sec:bridge-rh} 的 Abel--Weil 圆盘协议互补的“无参圆盘化”接口：将经典完成化函数
$$
\Xi_\zeta(s):=\tfrac12 s(s-1)\pi^{-s/2}\Gamma\!\left(\tfrac{s}{2}\right)\zeta(s)
$$
通过 Cayley 变换搬运为开单位圆盘内的解析函数，从而把经典 RH 等价地重述为“盘内零点禁入”与“一族 Jensen 缺陷恒零”。该重述的作用在于：它把无穷高度问题转写为单位圆盘内零点数据的否定性规格，并将零点信息进一步压缩为可审计的有限圆周平均。

\subsection{Cayley 变换与圆盘化完成化函数}
记上半平面与单位圆盘分别为
$$
\mathbb{H}:=\{t\in\CC:\ \Im(t)>0\},\qquad
\mathbb{D}:=\{w\in\CC:\ |w|<1\}.
$$
定义 Cayley 变换
$$
\boxed{\;
w=\mathcal{C}(t):=\frac{1+\mathrm{i}t}{1-\mathrm{i}t},
\qquad
t=\mathcal{C}^{-1}(w):=-\mathrm{i}\,\frac{w-1}{w+1}\; }.
$$

\begin{proposition}[Cayley 紧化的圆盘几何]\label{prop:app-cayley-upperhalf-disk-rh}
\(\mathcal{C}\) 给出双全纯同构 \(\mathcal{C}:\mathbb{H}\to\mathbb{D}\)。并且当 \(t\in\RR\) 时有 \(|\mathcal{C}(t)|=1\)，当 \(t\to\pm\infty\) 时有 \(\mathcal{C}(t)\to-1\)。
\end{proposition}
\begin{proof}
直接计算
$$
|\mathcal{C}(t)|^2=\frac{1+|t|^2-2\Im(t)}{1+|t|^2+2\Im(t)}.
$$
故 \(\Im(t)>0\Rightarrow |\mathcal{C}(t)|<1\)，\(\Im(t)=0\Rightarrow |\mathcal{C}(t)|=1\)。逆映射由代数求解得到，并可验证其把 \(|w|<1\) 映到 \(\Im(t)>0\)。当 \(t\in\RR\) 时 \(|\mathcal{C}(t)|=1\)，而 \(t\to\pm\infty\) 时 \(\mathcal{C}(t)\to-1\)。证毕。
\end{proof}

\paragraph{圆盘化完成化函数}
令
$$
\boxed{\;
s(w):=\frac12+\mathrm{i}\,\mathcal{C}^{-1}(w)=\frac12+\frac{w-1}{w+1}\; },
\qquad (|w|<1),
$$
并定义圆盘内解析函数
$$
\boxed{\;
F_\zeta(w):=\Xi_\zeta\!\bigl(s(w)\bigr)\qquad(|w|<1)\; }.
$$
由于 \(\Xi_\zeta\) 为整函数且 \(s(w)\) 在 \(|w|<1\) 内全纯（唯一极点 \(w=-1\) 位于边界），故 \(F_\zeta\) 在开单位圆盘内解析。
当 \(w=e^{\mathrm{i}\theta}\in\TT\setminus\{-1\}\) 时可写 \(\mathcal{C}^{-1}(e^{\mathrm{i}\theta})=\tan(\theta/2)\in\RR\)，从而
$$
F_\zeta(e^{\mathrm{i}\theta})=\Xi_\zeta\!\left(\tfrac12+\mathrm{i}\tan(\theta/2)\right),
$$
即把临界线截面参数化为单位圆边界相位。

\begin{theorem}[RH 的圆盘等价式：盘内零点禁入]\label{thm:app-rh-iff-disk-zero-free-rh}
经典 RH 等价于
$$
\boxed{\ \forall\,w\in\mathbb{D},\quad F_\zeta(w)\neq 0\ }.
$$
\end{theorem}
\begin{proof}
(\(\Rightarrow\)) 反设 RH 不成立，则存在 \(\Xi_\zeta\) 的零点 \(\rho=\beta+\mathrm{i}\gamma\) 满足 \(\beta\neq\tfrac12\) 且 \(0<\beta<1\)。由对称性 \(\Xi_\zeta(s)=\Xi_\zeta(1-s)\)，可不失一般性取 \(\beta<\tfrac12\)。
令
$$
t_\rho:=\gamma-\mathrm{i}\left(\beta-\tfrac12\right),
\qquad
\Im(t_\rho)=\tfrac12-\beta>0,
$$
则 \(t_\rho\in\mathbb{H}\) 且 \(\tfrac12+\mathrm{i}t_\rho=\rho\)。由命题 \ref{prop:app-cayley-upperhalf-disk-rh}，取 \(w_\rho:=\mathcal{C}(t_\rho)\in\mathbb{D}\) 并有 \(\mathcal{C}^{-1}(w_\rho)=t_\rho\)。因此
$$
F_\zeta(w_\rho)=\Xi_\zeta\!\left(\tfrac12+\mathrm{i}\,\mathcal{C}^{-1}(w_\rho)\right)
=\Xi_\zeta\!\left(\tfrac12+\mathrm{i}t_\rho\right)
=\Xi_\zeta(\rho)=0,
$$
从而盘内出现零点。
\par
(\(\Leftarrow\)) 反之，若存在 \(w_0\in\mathbb{D}\) 使 \(F_\zeta(w_0)=0\)，令 \(t_0:=\mathcal{C}^{-1}(w_0)\in\mathbb{H}\) 并置 \(s_0:=\tfrac12+\mathrm{i}t_0\)。由于 \(\Im(t_0)>0\)，有 \(\Re(s_0)=\tfrac12-\Im(t_0)<\tfrac12\)，且 \(\Xi_\zeta(s_0)=F_\zeta(w_0)=0\)。由函数方程与共轭对称性立刻推出存在偏离临界线的零点，故 RH 不成立。证毕。
\end{proof}

\paragraph{离临界线偏移的圆盘半径压缩}
定理 \ref{thm:app-rh-iff-disk-zero-free-rh} 的证明把离临界线零点事件搬运为圆盘内零点事件；下述命题给出该搬运在圆盘半径上的完全闭式压缩律。

\begin{proposition}[离临界线偏移的圆盘半径闭式压缩]\label{prop:app-offcritical-radius-compression-rh}
设 \(\Xi_\zeta\) 存在零点 \(\rho=\beta+\mathrm{i}\gamma\) 满足 \(\beta\neq\tfrac12\) 且 \(0<\beta<1\)。
由对称性 \(\Xi_\zeta(s)=\Xi_\zeta(1-s)\)，不失一般性取 \(\beta<\tfrac12\)，并令
$$
\delta:=\tfrac12-\beta>0,
\qquad
t_{\rho}:=\gamma+\mathrm{i}\delta\in\mathbb{H},
\qquad
w_{\rho}:=\mathcal{C}(t_{\rho})\in\mathbb{D}.
$$
则 \(F_\zeta(w_{\rho})=0\)，并且
$$
\boxed{\;
|w_{\rho}|^{2}
=\frac{\gamma^{2}+(\delta-1)^{2}}{\gamma^{2}+(\delta+1)^{2}},
\qquad
1-|w_{\rho}|^{2}
=\frac{4\delta}{\gamma^{2}+(\delta+1)^{2}}\; }.
$$
\end{proposition}
\begin{proof}
由定理 \ref{thm:app-rh-iff-disk-zero-free-rh} 的构造，\(F_\zeta(w_{\rho})=\Xi_\zeta(\rho)=0\)。
模长闭式由命题 \ref{prop:app-cayley-upperhalf-disk-rh} 的恒等式
\(
|\mathcal{C}(t)|^{2}=\frac{1+|t|^{2}-2\Im(t)}{1+|t|^{2}+2\Im(t)}
\)
取 \(t=t_{\rho}=\gamma+\mathrm{i}\delta\) 并化简得到。证毕。
\end{proof}

\begin{corollary}[高高度压缩尺度]\label{cor:app-offcritical-radius-high-height-rh}
当 \(|\gamma|\to\infty\) 且 \(\delta\) 有界时，
$$
1-|w_{\rho}|^{2}\sim \frac{4\delta}{\gamma^{2}},
\qquad
1-|w_{\rho}|\sim \frac{2\delta}{\gamma^{2}}.
$$
因此在固定偏移 \(\delta\) 下，圆盘坐标对偏移的可分辨尺度为 \(\Theta(\gamma^{-2})\)：高度越大，径向压缩越强。
\end{corollary}
\begin{proof}
由命题 \ref{prop:app-offcritical-radius-compression-rh} 得
\(
1-|w_{\rho}|^{2}=\frac{4\delta}{\gamma^{2}+(\delta+1)^{2}}\sim\frac{4\delta}{\gamma^{2}}
\)。
再由恒等式 \(1-|w_{\rho}|=\frac{1-|w_{\rho}|^{2}}{1+|w_{\rho}|}\) 且 \(|w_{\rho}|\to 1\) 得第二式。证毕。
\end{proof}

\subsection{Jensen 缺陷：把盘内零点计数压缩为一族有限圆周平均}
\begin{definition}[Jensen 缺陷]\label{def:app-jensen-defect-rh}
令 \(F_\zeta\) 如上节定义。对任意 \(0<\varrho<1\) 定义 Jensen 缺陷
$$
\boxed{\;
D_\zeta(\varrho)
:=\frac{1}{2\pi}\int_{0}^{2\pi}\log\bigl|F_\zeta(\varrho e^{\mathrm{i}\theta})\bigr|\,\dd\theta
-\log|F_\zeta(0)|\; }.
$$
\end{definition}

\begin{proposition}[Jensen 公式与非负缺陷]\label{prop:app-jensen-formula-defect-rh}
设 \(0<\varrho<1\)。令 \(\{a_k\}\) 为 \(F_\zeta\) 在 \(|w|<\varrho\) 内的零点（按重数计）。
则 Jensen 公式给出
$$
\boxed{\;
D_\zeta(\varrho)
=\sum_{|a_k|<\varrho}\log\!\Bigl(\frac{\varrho}{|a_k|}\Bigr)\ \ge\ 0\; }.
$$
从而 \(D_\zeta(\varrho)=0\) 当且仅当 \(F_\zeta\) 在 \(|w|<\varrho\) 内无零点。
\end{proposition}
\begin{proof}
由于 \(F_\zeta\) 在 \(|w|<1\) 内解析且 \(F_\zeta(0)=\Xi_\zeta(-\tfrac12)\neq 0\)，可对任意 \(\varrho<1\) 在圆盘 \(|w|<\varrho\) 上应用 Jensen 公式。
求和式以及非负性为标准推论，参见 \cite{Duren1970TheoryHp}。证毕。
\end{proof}

\begin{proposition}[单点盘内零点诱导的 Jensen 缺陷刚性下界]\label{prop:app-jensen-single-zero-lower-bound-rh}
若 \(F_\zeta\) 在开单位圆盘内存在零点 \(a\) 满足 \(|a|<1\)，则对任意 \(\varrho\in(|a|,1)\) 有
$$
\boxed{\;
D_\zeta(\varrho)\ge \log\!\Bigl(\frac{\varrho}{|a|}\Bigr),
\qquad
\liminf_{\varrho\uparrow 1}D_\zeta(\varrho)\ge \log\!\Bigl(\frac{1}{|a|}\Bigr)>0\; }.
$$
\end{proposition}
\begin{proof}
取任意 \(\varrho\in(|a|,1)\)。由命题 \ref{prop:app-jensen-formula-defect-rh} 的 Jensen 求和式，
对任意零点项 \(|a_k|<\varrho\) 均有非负贡献 \(\log(\varrho/|a_k|)\)，特别地包含 \(a\) 的贡献 \(\log(\varrho/|a|)\)，故
$$
D_\zeta(\varrho)=\sum_{|a_k|<\varrho}\log(\varrho/|a_k|)\ge \log(\varrho/|a|).
$$
令 \(\varrho\uparrow 1\) 即得 \(\liminf\) 下界。证毕。
\end{proof}

\begin{corollary}[RH 的完全有限化判据：缺陷恒零]\label{cor:app-rh-iff-jensen-defect-zero-rh}
经典 RH 等价于
$$
\boxed{\ \forall\,\varrho\in(0,1),\quad D_\zeta(\varrho)=0\ }.
$$
\end{corollary}
\begin{proof}
由命题 \ref{prop:app-jensen-formula-defect-rh}，对每个固定 \(\varrho\)，有 \(D_\zeta(\varrho)=0\iff F_\zeta\) 在 \(|w|<\varrho\) 内无零点。
令 \(\varrho\uparrow 1\) 得 “对一切 \(\varrho\in(0,1)\) 缺陷恒零” 等价于 “\(F_\zeta\) 在 \(|w|<1\) 内无零点”，再与定理 \ref{thm:app-rh-iff-disk-zero-free-rh} 合并即得。证毕。
\end{proof}

\begin{theorem}[趋近边界的缺陷趋零判据]\label{thm:app-rh-iff-jensen-defect-vanishing-sequence-rh}
下列命题等价：
\begin{enumerate}
  \item[(i)] 经典 RH 成立；
  \item[(ii)] 存在一列 \(\varrho_n\uparrow 1\) 使 \(D_\zeta(\varrho_n)\to 0\)。
\end{enumerate}
\end{theorem}
\begin{proof}
(i)\(\Rightarrow\)(ii) 由推论 \ref{cor:app-rh-iff-jensen-defect-zero-rh} 立得：若 RH 成立，则 \(D_\zeta(\varrho)\equiv 0\)，从而任取 \(\varrho_n\uparrow 1\) 即有 \(D_\zeta(\varrho_n)\to 0\)。
\par
(ii)\(\Rightarrow\)(i) 反设 RH 不成立。由定理 \ref{thm:app-rh-iff-disk-zero-free-rh}，存在盘内零点 \(a\in\mathbb{D}\) 使 \(F_\zeta(a)=0\)。
则由命题 \ref{prop:app-jensen-single-zero-lower-bound-rh} 得
$$
\liminf_{\varrho\uparrow 1}D_\zeta(\varrho)\ge \log(1/|a|)>0,
$$
与 \(D_\zeta(\varrho_n)\to 0\) 矛盾。故盘内零点不存在，进而由定理 \ref{thm:app-rh-iff-disk-zero-free-rh} 得 RH 成立。证毕。
\end{proof}

\begin{corollary}[缺陷上界序列给出的零点排除证书]\label{cor:app-jensen-defect-upperbound-certificate-rh}
若存在 \(\varrho_n\uparrow 1\) 与 \(\varepsilon_n\downarrow 0\) 使 \(D_\zeta(\varrho_n)\le \varepsilon_n\)，则 \(F_\zeta\) 在 \(\mathbb{D}\) 内无零点，从而经典 RH 成立。
更定量地，对任意盘内零点 \(a\in\mathbb{D}\)，当 \(n\) 充分大使 \(|a|<\varrho_n\) 时必有
$$
|a|\ge \varrho_n e^{-\varepsilon_n}.
$$
\end{corollary}
\begin{proof}
反设存在盘内零点 \(a\in\mathbb{D}\)。当 \(n\) 充分大使 \(|a|<\varrho_n\) 时，由命题 \ref{prop:app-jensen-single-zero-lower-bound-rh} 得
$$
\log(\varrho_n/|a|)\le D_\zeta(\varrho_n)\le \varepsilon_n,
$$
从而 \(|a|\ge \varrho_n e^{-\varepsilon_n}\)。
令 \(n\to\infty\) 得 \(|a|\ge 1\)，与 \(a\in\mathbb{D}\) 矛盾；故盘内零点不存在。由定理 \ref{thm:app-rh-iff-disk-zero-free-rh} 得 RH 成立。证毕。
\end{proof}

